\section{Methodology}
\label{sec:method}

\subsection{Task Design}
\label{sec:task_design}

We construct \benchname, a benchmark of 22 ordering tasks organized along two axes: \emph{property category} and \emph{difficulty level}.

\para{Property categories.}
Each task asks the model to sort a list of items by a single property.
We define five categories:
\begin{itemize}[leftmargin=*,itemsep=0pt,topsep=0pt]
    \item \textbf{Syntactic} (2 tasks): Properties determined by surface form (\eg alphabetical order of country or animal names).
    \item \textbf{Well-known factual} (3 tasks): Properties that form famous, commonly-taught sequences (\eg planet distance from the Sun, continent area).
    \item \textbf{General knowledge} (6 tasks): Approximate quantitative properties that educated adults could roughly estimate (\eg animal speeds, country populations).
    \item \textbf{Specific factual} (7 tasks): Precise numeric properties requiring exact recall (\eg city elevations, river lengths, element atomic weights).
    \item \textbf{Temporal} (4 tasks): Chronological ordering of events or people (\eg U.S.\ president order, invention dates, celebrity birth dates).
\end{itemize}

\para{Difficulty levels.}
Within each category, we assign tasks a difficulty level---easy (5 tasks), medium (8 tasks), or hard (9 tasks)---based on how commonly the ordering is encountered in educational or general-knowledge contexts.

\para{Item lists.}
Each task contains 6--10 items with verified ground-truth orderings.
Items are chosen to span the range of the property (\eg cities at diverse elevations) and to include pairs with similar values that test fine-grained discrimination.
\Tabref{tab:tasks} lists representative tasks.

\begin{table}[t]
    \centering
    \caption{Representative tasks from \benchname. Each task specifies a property and a list of 6--10 items with verified ground-truth orderings.}
    \label{tab:tasks}
    \resizebox{\textwidth}{!}{%
    \begin{tabular}{@{}llll@{}}
        \toprule
        \textbf{Task ID} & \textbf{Property} & \textbf{Category} & \textbf{Sample Items} \\
        \midrule
        \texttt{alpha\_countries} & Alphabetical order & Syntactic & France, Germany, Brazil, Japan, \ldots \\
        \texttt{planet\_distance} & Distance from Sun & Well-known & Saturn, Earth, Mars, Mercury, \ldots \\
        \texttt{country\_population} & Population & General knowledge & India, Nigeria, Japan, Brazil, \ldots \\
        \texttt{city\_elevation} & Elevation (m) & Specific factual & Denver, La Paz, Amsterdam, \ldots \\
        \texttt{celebrity\_dob} & Date of birth & Temporal & Taylor Swift, Elvis Presley, \ldots \\
        \bottomrule
    \end{tabular}
    }
\end{table}

\subsection{Experimental Protocol}
\label{sec:protocol}

\para{Prompt design.}
We use a standardized prompt: \texttt{``Sort the following items by [property]. Return ONLY the sorted list.''}
This minimal prompt avoids chain-of-thought scaffolding, ensuring that the model must access the relevant knowledge and produce the ordering in a single pass.

\para{Positional-bias control.}
For each task, we create three random shuffles of the input list (seeded with a fixed random seed of 42) and report the mean metric across shuffles.
This controls for positional bias, which \citet{tang2023found} show can substantially affect listwise ranking.

\para{Models.}
We evaluate two models: \gptlarge and \gptsmall.
\gptlarge represents a frontier-scale model; \gptsmall represents an efficient, smaller model.
Both are queried via the OpenAI API with temperature 0 and a maximum output length of 1{,}000 tokens.

\subsection{Evaluation Metrics}
\label{sec:metrics}

We report three metrics:
\begin{itemize}[leftmargin=*,itemsep=0pt,topsep=0pt]
    \item \textbf{\kendalltau} (primary): Rank correlation coefficient ranging from $-1$ (reversed) to $+1$ (perfect). Counts the fraction of concordant item pairs.
    \item \textbf{Exact match}: Binary indicator of whether the predicted ordering is identical to ground truth.
    \item \textbf{\spearmanrho}: Rank correlation that weights larger displacements more heavily than \kendalltau.
\end{itemize}

For statistical comparisons, we use the Wilcoxon signed-rank test (paired model comparison), the Kruskal--Wallis test (between-category comparison), and report Spearman correlation for cross-model agreement.
